% Source: https://artofproblemsolving.com/community/c3h1573318_problems_that_i_am_stuck
% Stable identity: aops-topic-1573318-q02
Put $y=-b$ and $z=-c$. The conditions become
\[
a+y+z=2015,\qquad\dfrac{1}{a}+\dfrac{1}{y}+\dfrac{1}{z}=\dfrac{1}{65},
\]
and since $\dfrac {a}{c}=-\dfrac {a}{z}$, $\dfrac {a}{b}=-\dfrac {a}{y}$ and so on, while $\dfrac {b}{c}=\dfrac {y}{z}$ and $\dfrac {c}{b}=\dfrac {z}{y}$, the expression equals
\[
-\left(\dfrac {a}{z}+\dfrac {z}{a}+\dfrac {a}{y}+\dfrac {y}{a}+\dfrac {y}{z}+\dfrac {z}{y}\right).
\]
Expanding the product of the two given sums,
\[
(a+y+z)\left(\dfrac{1}{a}+\dfrac{1}{y}+\dfrac{1}{z}\right)=3+\left(\dfrac {a}{y}+\dfrac {y}{a}+\dfrac {a}{z}+\dfrac {z}{a}+\dfrac {y}{z}+\dfrac {z}{y}\right),
\]
and the left side is $2015\cdot\dfrac{1}{65}=31$. So the bracket equals $31-3=28$, and the required value is $\boxed{-28}$.

Such numbers exist: for instance $a=100$ and $b,c$ the two (negative) roots $b=-y$, $c=-z$ where $y,z$ are the roots of $t^2-1915t+\dfrac{1915\cdot1300}{7}=0$, which has positive discriminant.
